\documentclass[../main.tex]{subfiles}
\begin{document}
	\thispagestyle{plain}
	
	\noindent
	\begin{tabular*}{\textwidth}{@{\extracolsep{\fill}}cc}
		\begin{minipage}[t]{0.46\textwidth}
			\centering
			\textbf{TRƯỜNG ĐẠI HỌC CÔNG NGHỆ}\\
			\textbf{THÔNG TIN VÀ TRUYỀN THÔNG}\\
			\textbf{KHOA CÔNG NGHỆ THÔNG TIN}\\[-0.15cm]
			\rule{5.2cm}{0.4pt}
		\end{minipage}
		&
		\begin{minipage}[t]{0.46\textwidth}
			\centering
			\textbf{CỘNG HÒA XÃ HỘI CHỦ NGHĨA VIỆT NAM}\\
			\textbf{Độc lập -- Tự do -- Hạnh phúc}\\[-0.15cm]
			\rule{5.5cm}{0.4pt}
		\end{minipage}
	\end{tabular*}
	
	\vspace{1.25cm}
	
	\begin{center}
		{\fontsize{15pt}{18pt}\selectfont\bfseries ĐỀ CƯƠNG ĐỒ ÁN / KHÓA LUẬN TỐT NGHIỆP}
	\end{center}
	
	\vspace{0.75cm}
	
	\fontsize{12pt}{17pt}\selectfont
	\setlength{\parindent}{0pt}
	
	\textbf{1. Thông tin sinh viên nhận đề tài}
	
	\vspace{0.22cm}
	Họ và tên: Hạng A Tráng
	
	\vspace{0.15cm}
	\begin{tabular*}{\textwidth}{@{\extracolsep{\fill}}ll}
		Mã sinh viên: DTC229200001 & Lớp: CNTT K21H\\[0.15cm]
		Điện thoại: 0349366182 & Email: dtc229200001@ictu.edu.vn
	\end{tabular*}
	
	\vspace{0.35cm}
	\textbf{2. Thông tin về giảng viên hướng dẫn}
	
	\vspace{0.22cm}
	Họ và tên: Nguyễn Lan Oanh
	
	\vspace{0.15cm}
	Học hàm, học vị: Thạc sĩ
	
	\vspace{0.15cm}
	\begin{tabular*}{\textwidth}{@{\extracolsep{\fill}}ll}
		Khoa: Công nghệ thông tin & Bộ môn: Công nghệ phần mềm\\[0.15cm]
		Điện thoại: 0948135145 & Email: nloanh@ictu.edu.vn
	\end{tabular*}
	
	\vspace{0.35cm}
	\textbf{3. Đề tài:} ``Xây dựng ứng dụng học tiếng Anh cho trẻ em có tích hợp AI hỗ trợ luyện phát âm''
	
	\vspace{0.35cm}
	\textbf{4. Thời gian thực hiện:} 10 tuần, từ ngày 09/03/2026 đến ngày 22/05/2026
	
	\vspace{0.35cm}
	\textbf{5. Mục tiêu:} Xây dựng một ứng dụng di động trên nền tảng Android hỗ trợ trẻ em từ 6 đến 10 tuổi học tiếng Anh hiệu quả. Ứng dụng tập trung vào việc tạo môi trường tương tác sinh động, tích hợp trí tuệ nhân tạo (AI) để nhận diện và đánh giá giọng nói, giúp trẻ luyện phát âm chuẩn xác ngay từ nhỏ, từ đó nâng cao kỹ năng nghe và nói một cách tự nhiên.
	
	\vspace{0.35cm}
	\textbf{6. Yêu cầu về chuẩn năng lực}
	
	Giáo viên không sửa nội dung này. Cố định theo hướng đề tài do lãnh đạo chuyên môn xây dựng.
	
	\vspace{0.3cm}
	
	\begin{center}
		\small
		\renewcommand{\arraystretch}{1.35}
		\setlength{\tabcolsep}{4pt}
		\begin{longtable}{|>{\centering\arraybackslash}p{0.8cm}|
				>{\raggedright\arraybackslash}p{3.8cm}|
				>{\centering\arraybackslash}p{0.8cm}|
				>{\centering\arraybackslash}p{0.8cm}|
				>{\raggedright\arraybackslash}p{5.4cm}|
				>{\centering\arraybackslash}p{0.8cm}|}
			\hline
			\textbf{Mã CLO} & \centering\textbf{Nội dung CLO} & \textbf{\%} & \textbf{Mã PI} & \centering\textbf{Nội dung PI} & \textbf{\%}\tabularnewline
			\hline
			\multirow{3}{*}{1} & \multirow{3}{3.8cm}{Phân tích được yêu cầu phần mềm} & \multirow{3}{*}{20}
			& 1.1 & Mô tả được quy trình nghiệp vụ & 5\\ \cline{4-6}
			& & & 1.2 & Mô tả được yêu cầu chức năng và phi chức năng & 5\\ \cline{4-6}
			& & & 1.3 & Xây dựng được sơ đồ UC và đặc tả UC & 10\\ \hline
			
			\multirow{3}{*}{2} & \multirow{3}{3.8cm}{Thiết kế được giải pháp cho phần mềm} & \multirow{3}{*}{30}
			& 2.1 & Thiết kế được cấu trúc dữ liệu & 15\\ \cline{4-6}
			& & & 2.2 & Thiết kế được kiến trúc phần mềm & 10\\ \cline{4-6}
			& & & 2.3 & Thiết kế mockup được giao diện & 5\\ \hline
			
			\multirow{4}{*}{3} & \multirow{4}{3.8cm}{Vận dụng một công cụ phát triển để thi công phần mềm} & \multirow{4}{*}{40}
			& 3.1 & Thi công được phần dữ liệu & 5\\ \cline{4-6}
			& & & 3.2 & Thi công đủ các thành phần theo kiến trúc & 5\\ \cline{4-6}
			& & & 3.3 & Thi công được các UC / yêu cầu chức năng & 20\\ \cline{4-6}
			& & & 3.4 & Thi công được các yêu cầu phi chức năng & 10\\ \hline
			
			\multirow{2}{*}{4} & \multirow{2}{3.8cm}{Triển khai thử nghiệm được phần mềm} & \multirow{2}{*}{10}
			& 4.1 & Triển khai phần mềm ra môi trường thử nghiệm & 5\\ \cline{4-6}
			& & & 4.2 & Tạo được dữ liệu thử nghiệm phù hợp & 5\\ \hline
		\end{longtable}
		\normalsize
	\end{center}
	
	\fontsize{12pt}{17pt}\selectfont
	
	\textbf{7. Yêu cầu về sản phẩm}
	
	\vspace{0.2cm}
	\textbf{a. Mô tả tóm tắt yêu cầu về sản phẩm của đồ án/khoá luận tốt nghiệp}
	
	Ứng dụng Android học tiếng Anh cho bé là một nền tảng di động hỗ trợ trẻ em từ 6 đến 10 tuổi tiếp cận ngôn ngữ mới một cách tự nhiên và sinh động. Sản phẩm cho phép trẻ học từ vựng và các mẫu câu cơ bản thông qua hình ảnh trực quan, âm thanh mẫu chuẩn, đồng thời tích hợp trí tuệ nhân tạo (AI) để nhận diện giọng nói và đánh giá kết quả phát âm của trẻ ngay lập tức.
	
	Đối với các bé, ứng dụng cung cấp giao diện tương tác vui nhộn, hỗ trợ luyện nghe và thực hành nói với các phản hồi khích lệ từ hệ thống AI giúp trẻ cải thiện phát âm một cách chính xác. Đối với phụ huynh, ứng dụng cung cấp công cụ theo dõi tiến độ học tập, ghi nhận kết quả và các kỹ năng bé đã đạt được theo thời gian thực.
	
	Về mặt kỹ thuật, ứng dụng yêu cầu giao diện thân thiện với trẻ em (UX Kids), tốc độ nhận diện giọng nói nhanh, tích hợp mượt mà các dịch vụ AI và đảm bảo tính ổn định trên các dòng thiết bị Android phổ biến.
	
	\vspace{0.2cm}
	\textbf{b. Yêu cầu chi tiết}
	
	\begin{center}
		\small
		\renewcommand{\arraystretch}{1.25}
		\setlength{\tabcolsep}{5pt}
		\begin{longtable}{|>{\centering\arraybackslash}p{1.2cm}|
				>{\raggedright\arraybackslash}p{13.2cm}|}
			\hline
			\textbf{PI} & \centering\textbf{Yêu cầu tối thiểu}\tabularnewline
			\hline
			1.1 & Có mô tả quy trình nghiệp vụ của ứng dụng học tiếng Anh cho trẻ em: từ khi bé chọn chủ đề học, xem từ vựng, nghe phát âm mẫu, thực hành phát âm và hệ thống AI nhận diện giọng nói và trả kết quả. Có sơ đồ quy trình nghiệp vụ (Flowchart) cho quá trình học và luyện phát âm.\\ \hline
			1.2 & Có bảng yêu cầu chức năng chi tiết của hệ thống như: Quản lý hồ sơ bé, Học từ vựng theo chủ đề, Luyện phát âm với AI, Xem bảng thành tích, Theo dõi tiến độ học tập. Ngoài ra có các yêu cầu phi chức năng như: giao diện thân thiện với trẻ em, tốc độ phản hồi nhanh của AI và đảm bảo dữ liệu học tập được lưu trữ an toàn trên thiết bị.\\ \hline
			1.3 & Có sơ đồ Use Case tổng quát của hệ thống. Có tối thiểu 5 Use Case chính gồm: Tạo hồ sơ bé, Chọn chủ đề học, Học từ vựng, Luyện phát âm với AI, Xem bảng thành tích. Có đặc tả chi tiết cho các Use Case quan trọng.\\ \hline
			2.1 & Có sơ đồ thực thể -- liên kết (ERD) cho cơ sở dữ liệu của ứng dụng gồm các bảng chính: Users (Bé), Categories (Chủ đề), Vocabulary (Từ vựng), PracticeHistory (Lịch sử luyện tập), Achievements (Thành tích). Có mô tả thuộc tính và mối quan hệ giữa các bảng.\\ \hline
			2.2 & Có sơ đồ thiết kế kiến trúc hệ thống cho ứng dụng Android. Kiến trúc gồm các thành phần: Android App (Java), cơ sở dữ liệu SQLite lưu trữ dữ liệu học tập và tích hợp dịch vụ nhận diện giọng nói (AI Speech Recognition API).\\ \hline
			2.3 & Có các bản thiết kế giao diện (Mockup) bằng Figma cho các màn hình chính: Trang chủ, Danh sách chủ đề học, Màn hình học từ vựng, Màn hình luyện phát âm với AI và bảng thành tích.\\ \hline
			3.1 & Xây dựng được cơ sở dữ liệu SQLite trên thiết bị Android theo thiết kế. Tạo dữ liệu bài học ban đầu gồm từ vựng, hình ảnh và âm thanh mẫu để phục vụ quá trình thử nghiệm.\\ \hline
			3.2 & Thi công được ứng dụng Android bằng Android Studio với ngôn ngữ Java. Ứng dụng có khả năng lưu trữ dữ liệu cục bộ và tích hợp thành công chức năng nhận diện giọng nói.\\ \hline
			3.3 & Thi công được tối thiểu 5 Use Case chính theo thiết kế. Hệ thống cho phép bé thực hiện toàn bộ quy trình học từ vựng và luyện phát âm, sau đó nhận đánh giá từ hệ thống.\\ \hline
			3.4 & Có các chức năng đặc trưng của ứng dụng như: hiệu ứng âm thanh và hình ảnh khi bé phát âm đúng, hệ thống huy hiệu hoặc điểm thưởng, và thống kê các từ bé thường phát âm sai để luyện lại.\\ \hline
			4.1 & Triển khai ứng dụng trên môi trường Android Emulator hoặc điện thoại Android thật. Ứng dụng chạy ổn định và các chức năng chính hoạt động bình thường.\\ \hline
			4.2 & Tạo dữ liệu thử nghiệm gồm ít nhất 50 từ vựng tiếng Anh kèm hình ảnh và âm thanh phát âm mẫu để kiểm thử các chức năng học tập và nhận diện giọng nói của hệ thống.\\ \hline
		\end{longtable}
		\normalsize
	\end{center}
	
	\fontsize{12pt}{17pt}\selectfont
	
	\textbf{8. Yêu cầu về quyển báo cáo:}
	
	Giảng viên không sửa nội dung này. Quyển báo cáo phải đảm bảo:
	\begin{itemize}
		\item Do sinh viên tự xây dựng;
		\item Có cấu trúc nội dung phù hợp yêu cầu và đảm bảo tính logic;
		\item Tuân thủ yêu cầu về khuôn mẫu và định dạng của trường;
		\item Không có lỗi chính tả;
		\item Không sử dụng nội dung vi phạm bản quyền;
		\item Không vi phạm quy định chống đạo văn của trường.
	\end{itemize}
	
	Cấu trúc của quyển báo cáo:
	\begin{itemize}
		\item Chương 1. Cơ sở lý thuyết (Những công nghệ chính sử dụng trong đồ án);
		\item Chương 2. Khảo sát và phân tích hệ thống;
		\item Chương 3. Thiết kế hệ thống;
		\item Chương 4. Triển khai (Cách cài đặt, bảo vệ, hoạt động thử nghiệm, tài liệu hướng dẫn).
	\end{itemize}
	
	\textbf{9. Tiến độ thực hiện}
	
	Giảng viên điều chỉnh cho phù hợp với đề tài. Đồ án thực hiện trong 10 tuần.
	
	\begin{center}
		\small
		\renewcommand{\arraystretch}{1.22}
		\setlength{\tabcolsep}{4pt}
		\begin{longtable}{|>{\centering\arraybackslash}p{0.8cm}|
				>{\centering\arraybackslash}p{1.45cm}|
				>{\raggedright\arraybackslash}p{6.1cm}|
				>{\raggedright\arraybackslash}p{6.1cm}|}
			\hline
			\textbf{STT} & \textbf{Thời gian} & \centering\textbf{Công việc} & \centering\textbf{Kết quả dự kiến}\tabularnewline
			\hline
			1 & Tuần 1 & Khảo sát thực tế nhu cầu học tiếng Anh của trẻ em từ 6--10 tuổi và nghiên cứu các ứng dụng học tiếng Anh phổ biến hiện nay. Phân tích cách trẻ tiếp cận bài học, cách hiển thị hình ảnh và âm thanh giúp trẻ dễ ghi nhớ từ vựng. & Danh sách các chủ đề học phù hợp với trẻ em (Animals, Fruits, Colors, Numbers...). Danh sách từ vựng cơ bản cho từng chủ đề. Sơ đồ quy trình học của ứng dụng: Chọn bài học $\rightarrow$ Xem từ vựng $\rightarrow$ Nghe phát âm $\rightarrow$ Luyện phát âm $\rightarrow$ AI đánh giá kết quả.\\ \hline
			2 & Tuần 2 & Phân tích yêu cầu hệ thống và xác định các chức năng cần thiết cho ứng dụng Android. Xây dựng bảng yêu cầu chức năng và phi chức năng của hệ thống. & Bảng yêu cầu chức năng: Quản lý hồ sơ bé, Học từ vựng theo chủ đề, Luyện phát âm bằng AI, Xem bảng thành tích, Lưu lịch sử học tập. Bảng yêu cầu phi chức năng: Giao diện dễ sử dụng cho trẻ em, tốc độ phản hồi nhanh, ứng dụng chạy ổn định trên Android.\\ \hline
			3 & Tuần 3 & Phân tích hệ thống và xây dựng sơ đồ Use Case cho ứng dụng. Mô tả chi tiết các Use Case chính của hệ thống. & Sơ đồ Use Case tổng quát của hệ thống. Đặc tả chi tiết các Use Case: Tạo hồ sơ bé, Chọn chủ đề học, Học từ vựng, Luyện phát âm với AI, Xem bảng thành tích.\\ \hline
			4 & Tuần 4 & Thiết kế cơ sở dữ liệu và kiến trúc hệ thống ứng dụng Android. Xác định các bảng dữ liệu cần thiết và mối quan hệ giữa chúng. & Sơ đồ ERD của hệ thống. Các bảng dữ liệu chính: Users, Categories, Vocabulary, PracticeHistory, Achievements. Sơ đồ kiến trúc hệ thống Android App kết nối với SQLite và AI Speech Recognition.\\ \hline
			5 & Tuần 5 & Thiết kế giao diện người dùng cho ứng dụng Android. Xây dựng các bản thiết kế giao diện (mockup) cho các màn hình chính của hệ thống. & Mockup giao diện bằng Figma. Các màn hình: Trang chủ, Danh sách chủ đề học, Màn hình học từ vựng, Màn hình luyện phát âm, Bảng thành tích.\\ \hline
			6 & Tuần 6 & Xây dựng cơ sở dữ liệu SQLite và chuẩn bị dữ liệu học tập ban đầu cho ứng dụng. Tạo các bảng dữ liệu và dữ liệu mẫu để phục vụ quá trình thử nghiệm. & Cơ sở dữ liệu SQLite hoàn chỉnh. Tạo dữ liệu mẫu gồm ít nhất 50 từ vựng kèm hình ảnh và âm thanh phát âm.\\ \hline
			7 & Tuần 7 & Thi công giao diện và các chức năng cơ bản của ứng dụng Android bằng Android Studio với ngôn ngữ Java. & Các Activity chính của ứng dụng: MainActivity, LessonActivity, PracticeActivity, ResultActivity. Ứng dụng hiển thị được danh sách chủ đề và từ vựng.\\ \hline
			8 & Tuần 8 & Tích hợp chức năng nhận diện giọng nói (AI Speech Recognition) để hỗ trợ luyện phát âm cho trẻ. Hoàn thiện các chức năng học tập chính của ứng dụng. & Chức năng ghi âm và nhận diện giọng nói hoạt động. Hệ thống đánh giá phát âm và hiển thị kết quả cho người học.\\ \hline
			9 & Tuần 9 & Kiểm thử toàn bộ ứng dụng, sửa lỗi giao diện, tối ưu chức năng và hoàn thiện trải nghiệm người dùng. & Ứng dụng chạy ổn định trên thiết bị Android. Các chức năng chính hoạt động đúng yêu cầu.\\ \hline
			10 & Tuần 10 & Hoàn thiện ứng dụng và hoàn thành báo cáo đồ án tốt nghiệp. Chuẩn bị tài liệu và nội dung trình bày bảo vệ đồ án. & File cài đặt ứng dụng Android (APK). Source code hoàn chỉnh. Báo cáo đồ án hoàn chỉnh.\\ \hline
		\end{longtable}
		\normalsize
	\end{center}
	
	\vspace{0.7cm}
	
	\begin{flushright}
		\textit{Thái Nguyên, ngày 06 tháng 03 năm 2026}
	\end{flushright}
	
	\vspace{0.4cm}
	
	\noindent
	\begin{tabular*}{\textwidth}{@{\extracolsep{\fill}}cc}
		\begin{minipage}{0.42\textwidth}
			\centering
			\textbf{LÃNH ĐẠO BỘ MÔN}
			
			\vspace{2.5cm}
		\end{minipage}
		&
		\begin{minipage}{0.42\textwidth}
			\centering
			\textbf{GIẢNG VIÊN HƯỚNG DẪN}
			
			\vspace{2.2cm}
			
			\textbf{Nguyễn Lan Oanh}
		\end{minipage}
	\end{tabular*}
	
\end{document}